% !PRINTED 24
\weberpage{24}
%<<CONT
\begingroup
\setlength{\abovedisplayskip}{5pt plus 2pt minus 1pt}
\setlength{\belowdisplayskip}{5pt plus 2pt minus 1pt}
\setlength{\abovedisplayshortskip}{3pt plus 1pt}
\setlength{\belowdisplayshortskip}{4pt plus 1pt}
\linespread{0.94}\selectfont
gleich dem höchsten der Grade der Summanden und kann sich nur in dem
besonderen Falle erniedrigen, wenn der höchste Grad in mehreren
Summanden vorkommt und die Summe der Coëfficienten der höchsten
Potenzen gleich Null ist.

Bei der Multiplication zweier oder mehrerer ganzer rationaler
Functionen entsteht eine Function, deren Grad gleich der Summe der
Grade der Factoren ist.

Um für das Product das Bildungsgesetz der Coëfficienten zu übersehen,
setzen wir
\begin{equation}
\begin{aligned}
 A(x)&=a_0x^m+a_1x^{m-1}+a_2x^{m-2}+\cdots,\\
 B(x)&=b_0x^n+b_1x^{n-1}+b_2x^{n-2}+\cdots,
\end{aligned}
\tag{2}
\end{equation}
\begin{equation}
 A(x)B(x)=C(x)=c_0x^{m+n}+c_1x^{m+n-1}+c_2x^{m+n-2}+\cdots
\tag{3}
\end{equation}
und erhalten
\[
\begin{aligned}
 c_0&=a_0b_0,\\
 c_1&=a_0b_1+a_1b_0,\\
 c_2&=a_0b_2+a_1b_1+a_2b_0,\\
 &\cdots,
\end{aligned}
\]
oder in allgemeinen Zeichen, wenn $\nu$ eine der Zahlen
$0,1,2,\ldots,m+n$ bedeutet:
\begin{equation}
 c_\nu=a_0b_\nu+a_1b_{\nu-1}+a_2b_{\nu-2}+\cdots
       +a_{\nu-1}b_1+a_\nu b_0,
\tag{4}
\end{equation}
worin alle $a$, deren Index grösser als $m$, und alle $b$, deren Index
grösser als $n$ ist, gleich Null zu setzen sind. Denn $c_\nu$ ist der
Coëfficient von $x^{m+n-\nu}$, und es entsteht also
$c_\nu x^{m+n-\nu}$ durch Multiplication aller Glieder von der Form
$a_\mu x^{m-\mu}$ mit allen Gliedern der Form
$b_{\nu-\mu}x^{n-\nu+\mu}$ und darauf folgende Addition. Sind $a_0$
und $b_0=1$, so ist auch $c_0=1$.% [source-original wording] The print applies “=1” only after b_0; the intended joint condition is preserved without silently rewriting the German.

\gesperrt{Hat in allen Factoren eines solchen Products die höchste
Potenz von $x$ den Coëfficienten $1$, so ist auch im Product der
Coëfficient der höchsten Potenz $=1$.}

Aus diesen Vorschriften für die Rechnung mit ganzen Functionen folgt,
dass die Regeln des Rechnens, die sich in Formeln wie $ab=ba$,
$(ab)c=a(bc)$, $(a+b)c=ac+bc$ und ähnlichen aussprechen, auch wenn
$a,b,c$ ganze Functionen von $x$ sind, gelten, und zwar in dem Sinne,
dass ganze Functionen nur dann als gleich gelten, wenn gleich hohe
Potenzen gleiche Coëfficienten haben.

Unter ganzen Functionen mehrerer Veränderlichen verstehen wir Summen
von Producten von ganzen, positiven Potenzen der Veränderlichen
$x,y,z\ldots$ mit irgend welchen Coëfficienten, die als constant
gelten. Man kann sich diese Functionen entstanden
\endgroup
%CONT>>
